\documentclass[11pt]{article}
\usepackage{amsmath,amssymb,amsthm}
\usepackage[margin=1.1in]{geometry}

\title{Minimization Algorithms: From Newton--Raphson to Adam}
\author{Yuval Lavie}
\date{\today}

\begin{document}
\maketitle

\noindent
A lineage of update rules for $\min_\theta f(\theta)$, each fixing a
weakness of its predecessor. All are implemented from scratch and raced on
the Rosenbrock valley in \texttt{optimizers.py}. Throughout, $g_k =
\nabla f(\theta_k)$ and $\eta$ is the step size (learning rate).

\section{Second order: Newton--Raphson}

Newton's method finds roots of an equation; \emph{optimization is
root-finding on the gradient} ($\nabla f = 0$). Modeling $f$ locally by
its second-order Taylor expansion and jumping to that model's minimum:
\begin{equation}
  \theta_{k+1} \;=\; \theta_k - H_k^{-1} g_k ,
  \qquad H_k = \nabla^2 f(\theta_k).
  \label{eq:newton}
\end{equation}
Near a minimum with $H \succ 0$, convergence is \emph{quadratic} — the
number of correct digits doubles per step. The costs: computing and
solving a $d\times d$ Hessian per step ($O(d^3)$ — hopeless for large
models), and far from the optimum $H_k$ may fail to be positive definite,
sending \eqref{eq:newton} uphill (the script damps it when that happens).
Everything after this section is an attempt to keep some of Newton's
geometry while paying first-order prices.

\section{First order: gradient descent}

\begin{equation}
  \theta_{k+1} \;=\; \theta_k - \eta\, g_k .
  \label{eq:gd}
\end{equation}
For $L$-smooth convex $f$ and $\eta \le 1/L$: $f(\theta_k)-f^\ast =
O(1/k)$; add $\mu$-strong convexity and the rate is linear (geometric)
with contraction governed by the condition number $\kappa = L/\mu$. That
$\kappa$-dependence is GD's weakness: on an ill-conditioned valley
(Rosenbrock) the largest stable $\eta$ is set by the steep direction,
leaving the shallow direction to crawl.

\paragraph{SGD.} Replace $g_k$ by an unbiased estimate from one sample or
a mini-batch (as in \texttt{bernoulli/}, \texttt{gaussian/},
\texttt{linear-regression/}): cost per step drops from $O(n)$ to $O(1)$
samples; a constant step leaves a noise ball, handled by averaging
(Polyak--Ruppert) or decaying steps. All methods below take stochastic
gradients in practice.

\section{Momentum: heavy ball and Nesterov}

\textbf{Heavy ball (Polyak).} Accumulate a velocity:
\begin{equation}
  v_{k+1} = \beta v_k + g_k,
  \qquad
  \theta_{k+1} = \theta_k - \eta\, v_{k+1},
  \qquad \beta \in [0,1).
  \label{eq:hb}
\end{equation}
Along a persistent downhill direction the geometric series boosts the
effective step by up to $1/(1-\beta)$; across the valley, alternating
gradients cancel in $v$. Physically: a ball with inertia rolling on the
loss surface with friction $1-\beta$.

\textbf{Nesterov accelerated gradient.} Evaluate the gradient at the
\emph{look-ahead} point:
\begin{equation}
  v_{k+1} = \beta v_k + \nabla f\bigl(\theta_k - \eta\beta v_k\bigr),
  \qquad
  \theta_{k+1} = \theta_k - \eta\, v_{k+1}.
  \label{eq:nesterov}
\end{equation}
The correction ``looks where the momentum is taking you before deciding.''
On smooth convex problems it achieves $f(\theta_k)-f^\ast = O(1/k^2)$ —
provably optimal for first-order methods, versus GD's $O(1/k)$.

\section{Quasi-Newton: BFGS and L-BFGS}

Estimate curvature from gradient differences instead of computing $H$.
With $s_k = \theta_{k+1}-\theta_k$ and $y_k = g_{k+1}-g_k$, BFGS maintains
an inverse-Hessian estimate $B_k$ satisfying the \emph{secant condition}
$B_{k+1} y_k = s_k$:
\begin{equation}
  B_{k+1} = \bigl(I - \rho_k s_k y_k^\top\bigr) B_k
            \bigl(I - \rho_k y_k s_k^\top\bigr)
            + \rho_k s_k s_k^\top,
  \qquad \rho_k = \frac{1}{y_k^\top s_k},
  \label{eq:bfgs}
\end{equation}
stepping $\theta_{k+1} = \theta_k - \alpha_k B_k g_k$ with a line search
($y_k^\top s_k > 0$ keeps $B \succ 0$). Superlinear convergence at
first-order cost per step. \emph{L-BFGS} stores only the last $m$ pairs
$(s,y)$ and applies $B_k$ implicitly — $O(md)$ memory, the standard for
smooth medium-scale problems. (Quasi-Newton dislikes stochastic gradients:
noisy $y_k$ corrupts the curvature estimate, one reason deep learning went
the adaptive route instead.)

\section{Adaptive per-coordinate steps: AdaGrad, RMSProp, Adam}

One global $\eta$ misfits problems whose coordinates have very different
scales. Give each coordinate its own step from its gradient history
(all operations below are element-wise).

\textbf{AdaGrad:}
\begin{equation}
  s_{k+1} = s_k + g_k^2,
  \qquad
  \theta_{k+1} = \theta_k - \eta\, \frac{g_k}{\sqrt{s_{k+1}}+\varepsilon}.
  \label{eq:adagrad}
\end{equation}
Frequently/steeply updated coordinates get shrinking steps. Flaw: $s$ only
grows, so all steps decay to zero eventually.

\textbf{RMSProp} forgets exponentially, fixing the decay:
\begin{equation}
  s_{k+1} = \rho\, s_k + (1-\rho)\, g_k^2,
  \qquad
  \theta_{k+1} = \theta_k - \eta\, \frac{g_k}{\sqrt{s_{k+1}}+\varepsilon}.
  \label{eq:rmsprop}
\end{equation}

\textbf{Adam} = RMSProp + momentum + bias correction:
\begin{equation}
  \begin{aligned}
    m_{k+1} &= \beta_1 m_k + (1-\beta_1)\, g_k,
    &\hat m &= m_{k+1}/(1-\beta_1^{\,k+1}),\\
    v_{k+1} &= \beta_2 v_k + (1-\beta_2)\, g_k^2,
    &\hat v &= v_{k+1}/(1-\beta_2^{\,k+1}),
  \end{aligned}
  \qquad
  \theta_{k+1} = \theta_k - \eta\, \frac{\hat m}{\sqrt{\hat v}+\varepsilon}.
  \label{eq:adam}
\end{equation}
The $(1-\beta^{k})$ corrections debias the zero-initialized moment
estimates in early iterations. Defaults $\beta_1=0.9$, $\beta_2=0.999$
have proven remarkably robust — Adam is the de-facto default of deep
learning. \emph{AdamW} decouples weight decay from the adaptive gradient
step ($\theta \leftarrow \theta - \eta\lambda\theta$ applied separately),
the variant modern training recipes actually use.

\section{What the race shows}

On Rosenbrock from $(-1.5, 2)$ (see the script's log-scale figure):
plain GD crawls along the valley floor; momentum and Nesterov accelerate
it substantially; the adaptive family navigates the bad scaling without
hand-tuned per-direction steps; and the curvature methods --- Newton and
BFGS --- drop to machine precision in tens of steps, which is what
quadratic/superlinear convergence looks like on a plot. No single winner:
the ranking would change with dimension, stochasticity, and cost per
step --- which is precisely why the zoo exists.

\end{document}
